\section{Method}

We consider a dataset of cells $\{x_i\}_{i=1}^N$ with $B \geq 1$ batches, where $x_i \in \mathbb{R}^G$ denotes gene expression and $b_i \in \{1, \dots, B\}$ is the batch label. The model consists of a batch-conditional encoder mapping cells to a shared latent space, a set of learnable prototypes representing metacells, and a batch-conditional decoder reconstructing gene expression from prototypes, so that batch-specific variation is handled at decoding rather than absorbed into shared prototypes.

\subsection{Model Components}

\paragraph{Input.}
We assume gene expression profiles $x_i \in \mathbb{R}^G$ with batch labels $b_i$, and a precomputed cell--cell similarity graph $P \in \mathbb{R}^{N \times N}$. By default, $P$ is an adaptive Gaussian (RBF) kernel on PCA embeddings~\citep{seacells}, which corrects for local density variations to ensure rare cell types form well-connected communities.
If the affinity is computed per batch, the resulting graphs can simply be concatenated to form a block-diagonal $P$; even when computed globally across all cells, batch effects make cross-batch similarities negligible, so $P$ is approximately block-diagonal in either case (see Appendix~\ref{app:graph_structure}).
Since the affinity is constructed in PCA space, it reflects both biological variation and technical batch effects; disentangling these signals is delegated to the batch-conditional encoder during training.
More broadly, the affinity can encode any biologically meaningful similarity---including spatial context---allowing the model to capture microenvironmental structure while remaining grounded in gene expression via the reconstruction objective (see Appendix~\ref{app:spatial_affinity}).

\paragraph{Encoder.}
A conditional encoder maps each cell to latent space:
\[
z_i = f_\theta(x_i, b_i) \in \mathbb{R}^d.
\]

\paragraph{Prototypes.}
We introduce $K$ learnable prototypes $\{c_k\}_{k=1}^K$, $c_k \in \mathbb{R}^d$, representing metacells in the latent space.

\paragraph{Soft Assignment.}
Each cell is represented as a soft mixture of prototypes:
\[
s_{ik} =
\frac{\exp\left( \frac{z_i^\top c_k}{\tau} \right)}
{\sum_{k'} \exp\left( \frac{z_i^\top c_{k'}}{\tau} \right)},
\]
where $\tau > 0$ is a temperature parameter controlling assignment sharpness, and both $z_i$ and $c_k$ are $\ell_2$-normalized (see Appendix~\ref{app:training}). Rather than treating $\tau$ as a fixed hyperparameter, we calibrate it from data after pretraining and prototype initialization (see Appendix~\ref{app:calibration}).

The assignment vector $s_i \in \mathbb{R}^K$ encodes the soft membership of cell $i$ across all $K$ prototypes. This representation is used to preserve similarity structure by mapping highly connected cells to similar assignment vectors.

\paragraph{Community Structure Preservation.}
To preserve local similarity structure, we define similarity between cells based on prototype assignments:
\[
q_{ij} = s_i^\top s_j.
\]

Since $s_i$ is a probability vector (softmax outputs are non-negative and sum to one), $q_{ij} \in [0,1]$, making it a natural similarity for binary cross-entropy.

We match this similarity to the input graph using a binary cross-entropy objective:
\[
\mathcal{L}_{\text{community}}
=
-\sum_{i,j}
\Big[
P_{ij}\log \tilde q_{ij}
+
(1-P_{ij})\log(1-\tilde q_{ij})
\Big],
\]
where $\tilde q_{ij}$ clips $q_{ij}$ away from 0 and 1 to avoid $\log(0)$.

This objective encourages cells strongly connected in the affinity graph to share prototype assignments. In practice, we optimize stochastically, sampling positive edges from $P_{ij}$ and negative pairs within each batch; since both operate within individual datasets, no assumption is made about cross-batch similarity or dissimilarity, and cross-batch alignment is delegated to the encoder via the reconstruction objective.

\paragraph{Decoder and Reconstruction.}
Each prototype is decoded into gene expression space:
\[
\hat{x}_k^{(b)} = g_\phi(c_k,\, b),
\]
and each cell is reconstructed as:
\[
\tilde{x}_i = \sum_{k=1}^K s_{ik}\, \hat{x}_k^{(b_i)}.
\]

We optimize:
\[
\mathcal{L}_{\text{rec}} = \frac{1}{N} \sum_i \|x_i - \tilde{x}_i\|_2^2.
\]

In practice, we stop the gradient through $s_{ik}$ when optimizing $\mathcal{L}_{\text{rec}}$, so the encoder receives no gradient from $\mathcal{L}_{\text{rec}}$, preventing reconstruction from conflicting with the affinity-based objectives (see Appendix~\ref{app:training}).

\paragraph{Prototype Regularization.}

\textbf{Batch-aware normalized association.}

Let $S \in \mathbb{R}^{N \times K}$ denote the assignment matrix.
For each batch $b$, we define a batch-specific edge matrix:
\[
W^{(b)}_{ij} =
\begin{cases}
P_{ij} & \text{if } b_i = b \;\text{or}\; b_j = b, \\0 & \text{otherwise}.
\end{cases}
\]

Let $d^{(b)}$ denote the corresponding degree vector, where
\[
d^{(b)}_i = \sum_j W^{(b)}_{ij}.
\]

We then compute:
\[
A^{(b)} = S^\top W^{(b)} S, \quad
\mathrm{vol}^{(b)} = S^\top d^{(b)},
\]
\[
M^{(b)}_{kj} =
\frac{A^{(b)}_{kj}}
{\sqrt{\mathrm{vol}^{(b)}_k \, \mathrm{vol}^{(b)}_j}}.
\]

Diagonal terms $M^{(b)}_{kk} \in [0,1]$ measure within-batch prototype compactness (1 = tight, well-connected community); off-diagonal terms $M^{(b)}_{kj}$ measure inter-prototype connectivity (0 = distinct, non-overlapping).

To aggregate across batches, we use:
\[
M_{kj} = \max_b M^{(b)}_{kj}.
\]

For diagonal terms, a prototype is considered compact if it forms a tight community in at least one batch; for off-diagonal terms, minimizing the maximum ensures no batch exhibits strong mixing between prototypes.

The loss is:
\[
\mathcal{L}_{\text{nassoc}} =
\frac{1}{K} \sum_k (M_{kk} - 1)^2
+
\alpha \cdot \frac{1}{K(K-1)} \sum_{k \neq j} M_{kj}^2.
\]

\textbf{Prototype usage.}
To prevent inactive prototypes, we penalize those with low maximum assignment weight.
Rather than using a noisy per-mini-batch estimate, we maintain a scalar exponential moving average (EMA) $u_k \in \mathbb{R}$ across mini-batches with momentum $m \in (0,1)$:
\[
u_k \leftarrow m\, u_k + (1-m)\, \max_i s_{ik},
\]
and penalize underused prototypes:
\[
\mathcal{L}_{\text{usage}} = \frac{1}{K} \sum_k -\log(u_k).
\]

Together with the normalized association objective, this ensures that each prototype represents a distinct and meaningful community.

\paragraph{Total Objective.}
\[
\mathcal{L} =
\mathcal{L}_{\text{community}}
+
\lambda \mathcal{L}_{\text{rec}}
+
\lambda_n \mathcal{L}_{\text{nassoc}}
+
\lambda_u \mathcal{L}_{\text{usage}}.
\]

\paragraph{Training Procedure.}

\textbf{Stage 1: Pretraining.}
We pretrain the encoder and decoder using scPoli~\citep{scpoli}, a cVAE-based architecture for single-cell data that encodes batch identity via trainable batch embeddings rather than fixed one-hot vectors, enabling effective batch effect correction.
We optimize a per-cell MSE reconstruction loss on log-normalized gene expression:
\[
\mathcal{L}_{\text{cVAE}} =
\frac{1}{N}\sum_i \|x_i - g_\phi(z_i, b_i)\|_2^2,
\]
where $z_i = f_\theta(x_i, b_i)$. This stage trains the encoder and batch embeddings before prototypes are introduced (see Appendix~\ref{app:training} for architecture and hyperparameter details).

\textbf{Stage 2: Prototype initialization and learning.}
Prototypes are initialized via waypoint initialization~\citep{seacells}: $K$ cells selected by greedy MaxMin in diffusion map space, with pretrained encoder embeddings as initial prototype vectors. The temperature $\tau$ is calibrated before joint optimization (Appendix~\ref{app:calibration}), and the KL term is omitted since pretraining already regularizes the latent space.

This two-stage procedure stabilizes training and yields prototypes that preserve within-batch structure, align shared cell states across batches, and represent compact, distinct communities.
Full hyperparameter details, loss weights, and guidance on selecting $K$ are provided in Appendix~\ref{app:training}.